\documentclass[12pt]{ximera}
\title{Exam 2}
\begin{document}
\begin{abstract}
Exam 2, covering Sections 7.5--8.2: properties of the counting formulas, conditional
probability and independence from survey data, choosing the right counting technique,
and a Bayes' Theorem calculation for a veterinary test.
\end{abstract}
\maketitle

\section*{Exam 2}

Give probabilities as exact fractions, decimals, or percentages; do not round unless
specifically asked to do so.

Formulas that may be useful:
\[
P(E)=\frac{n(E)}{n(S)}, \qquad P(E\mid F)=\frac{P(E\cap F)}{P(F)},
\]
\[
P(N,k)=\frac{N!}{(N-k)!}, \qquad C(N,k)=\frac{N!}{k!\,(N-k)!}.
\]

\begin{problem}
Assess whether each of the following claims is true or false.

\begin{prompt}
\textbf{(a)} Claim: the formulas for permutations and combinations assume that
elements CANNOT be selected more than once.

\begin{multipleChoice}
\choice[correct]{True}
\choice{False}
\end{multipleChoice}

\begin{solution}
True. Both formulas count selections \emph{without} replacement --- that is exactly
why the factors run $N, N-1, N-2,\ldots$ instead of staying at $N$. When repetition
\emph{is} allowed, you use the multiplication principle instead, as in $4^7$ or
$13^5$.
\end{solution}
\end{prompt}

\begin{prompt}
\textbf{(b)} Claim: the multiplication principle does NOT work in situations where
order matters.

\begin{multipleChoice}
\choice{True}
\choice[correct]{False}
\end{multipleChoice}

\begin{solution}
False --- it is the reverse. The multiplication principle works particularly well
when order matters; the permutation formula is derived from it. For instance,
choosing first, second, and third place among $20$ students is
$20\cdot 19\cdot 18$, a direct application.

Where the multiplication principle needs correcting is when order does \emph{not}
matter, since then it counts each unordered selection several times.
\end{solution}
\end{prompt}

\begin{prompt}
\textbf{(c)} Claim: the formulas for permutations and combinations ALWAYS compute to
a whole number.

\begin{multipleChoice}
\choice[correct]{True}
\choice{False}
\end{multipleChoice}

\begin{solution}
True. Both formulas count a number of things --- arrangements or subsets --- and you
cannot have a fractional number of subsets, so the result must be a whole number.

This is worth remembering as an error check: $C(N,k)$ contains a division by
$k!\,(N-k)!$, and if your arithmetic ever leaves you with a fraction or decimal, you
have made a mistake somewhere.
\end{solution}
\end{prompt}
\end{problem}

\begin{problem}
Suppose that in a survey of $70$ students, $50$ report they like wallabies and $35$
report they like axolotls; $20$ of the students report they like \emph{both}.

\begin{prompt}
\textbf{(a)} Among the students who report that they like wallabies, what is the
probability that they also like axolotls?

Give your answer as a fraction in lowest terms: $\answer{2/5}$

\begin{hint}
This is $P(A\mid W)$ --- restrict attention to the $50$ wallaby-likers and ask how
many of them also like axolotls.
\end{hint}

\begin{solution}
Restricting to the $50$ students who like wallabies, $20$ of them also like axolotls:
\[
P(A\mid W)=\frac{n(A\cap W)}{n(W)}=\frac{20}{50}=\frac25=0.4.
\]
\end{solution}
\end{prompt}

\begin{prompt}
\textbf{(b)} Recall that if $P(E\mid F)=P(E)$ then the events $E$ and $F$ are said to
be \emph{independent}. Based on this data, are liking wallabies and liking axolotls
independent events?

\begin{multipleChoice}
\choice{Yes}
\choice[correct]{No}
\end{multipleChoice}

\begin{solution}
No. From part (a), $P(A\mid W)=0.4$. But overall,
\[
P(A)=\frac{35}{70}=0.5.
\]
Since $0.4\neq 0.5$, the events are not independent --- knowing a student likes
wallabies makes it somewhat \emph{less} likely that they like axolotls.
\end{solution}
\end{prompt}
\end{problem}

\begin{problem}
Determine whether the following should be counted using the multiplication principle,
permutations, or combinations, and give the count.

\begin{prompt}
\textbf{(a)} You are ranking your top $5$ favorite Pok\'emon from among the $100$
introduced in second-gen to compare with friends. How many different top-$5$ rankings
are possible?

\begin{multipleChoice}
\choice[correct]{Permutations}
\choice{Combinations}
\choice{Multiplication principle with repetition}
\end{multipleChoice}

The count is $\answer{9034502400}$.

\begin{solution}
A ranking is ordered, and no Pok\'emon appears twice, so this is a permutation:
\[
P(100,5)=100\cdot 99\cdot 98\cdot 97\cdot 96 = 9{,}034{,}502{,}400.
\]
\end{solution}
\end{prompt}

\begin{prompt}
\textbf{(b)} Each week, Saturday Night Live comedians propose $40$ sketches, and by
the end of the week they select $10$ of the sketches to perform. How many different
ways could they choose which $10$ to perform?

\begin{multipleChoice}
\choice{Permutations}
\choice[correct]{Combinations}
\choice{Multiplication principle with repetition}
\end{multipleChoice}

The count is $\answer{847660528}$.

\begin{solution}
The question asks only \emph{which} ten are chosen, not in what order:
\[
C(40,10)=\frac{40!}{10!\,30!}=847{,}660{,}528.
\]
\end{solution}
\end{prompt}
\end{problem}

\begin{problem}
Calculate the total number of possible outcomes for each of the following.

\begin{prompt}
\textbf{(a)} You have $3$ homework assignments to complete and want to allot a
particular day of the week (Sunday--Saturday) to complete each one, \emph{without}
scheduling multiple assignments for the same day. How many different ways could you
schedule these assignments?

$\answer{210}$

\begin{solution}
The assignments are distinct and no day is reused, so this is a permutation of $3$
days chosen from $7$:
\[ P(7,3)=7\cdot 6\cdot 5 = 210. \]
\end{solution}
\end{prompt}

\begin{prompt}
\textbf{(b)} A group of $5$ friends are building D\&D characters. How many ways can
they all select their character's class if there are $13$ possible classes and
multiple players \emph{may} choose the same class?

$\answer{371293}$

\begin{hint}
Repeats are allowed here, so this is the multiplication principle. Which number is
the base and which is the exponent?
\end{hint}

\begin{solution}
Each of the $5$ friends independently picks one of $13$ classes, and repeats are
allowed:
\[ 13^5 = 371{,}293. \]
Note it is $13^5$, not $5^{13}$: there are $5$ decisions to make, each with $13$
options.
\end{solution}
\end{prompt}
\end{problem}

\begin{problem}
Suppose the probability of a dog in a particular city having heartworm disease (H) is
estimated to be approximately $0.01$. One test for heartworm \emph{accurately}
identifies positive cases with a probability of $0.7$ and \emph{accurately} identifies
negative cases with a probability of $0.9$.

\begin{prompt}
\textbf{(a, first branches)} $P(H)=\answer[tolerance=0.0001]{0.01}$ \qquad
$P(\text{no H})=\answer[tolerance=0.0001]{0.99}$

\begin{solution}
The disease rate is given as $0.01$, and its complement is $1-0.01=0.99$.
\end{solution}
\end{prompt}

\begin{prompt}
\textbf{(a, second branches)} Given the dog has heartworm: \\
$P(\text{pos}\mid H)=\answer[tolerance=0.0001]{0.7}$ \qquad
$P(\text{neg}\mid H)=\answer[tolerance=0.0001]{0.3}$

\smallskip
Given the dog does not: \\
$P(\text{pos}\mid \text{no H})=\answer[tolerance=0.0001]{0.1}$ \qquad
$P(\text{neg}\mid \text{no H})=\answer[tolerance=0.0001]{0.9}$

\begin{hint}
``Accurately identifies positive cases with probability $0.7$'' is the sensitivity,
$P(\text{pos}\mid H)$. The other two branches in each pair are complements.
\end{hint}

\begin{solution}
Accurately identifying positive cases means $P(\text{pos}\mid H)=0.7$, so the
false-negative branch is $P(\text{neg}\mid H)=1-0.7=0.3$.

Accurately identifying negative cases means $P(\text{neg}\mid\text{no H})=0.9$, so the
false-positive branch is $P(\text{pos}\mid\text{no H})=1-0.9=0.1$.
\end{solution}
\end{prompt}

\begin{prompt}
\textbf{(a, outcomes)} Multiply along each path. Round to four decimal places.

$P(H\cap\text{pos})=\answer[tolerance=0.00005]{0.007}$ \qquad
$P(H\cap\text{neg})=\answer[tolerance=0.00005]{0.003}$

\smallskip
$P(\text{no H}\cap\text{pos})=\answer[tolerance=0.00005]{0.099}$ \qquad
$P(\text{no H}\cap\text{neg})=\answer[tolerance=0.00005]{0.891}$

\begin{solution}
\[
\begin{array}{ll}
P(H\cap\text{pos})=(0.01)(0.7)=0.0070 &
P(H\cap\text{neg})=(0.01)(0.3)=0.0030\\[2pt]
P(\text{no H}\cap\text{pos})=(0.99)(0.1)=0.0990 &
P(\text{no H}\cap\text{neg})=(0.99)(0.9)=0.8910
\end{array}
\]
These four add to $1.0000$. \checkmark
\end{solution}
\end{prompt}

\begin{prompt}
\textbf{(b)} If a dog from this city tests positive for heartworm, what is the
probability that the dog does in fact have heartworm disease? That is, what is
$P(H\mid \text{pos})$?

Give your final answer as a percent, rounded to the nearest percent:
$\answer{7}\%$

\begin{hint}
Add the two paths that end in a positive test to get $P(\text{pos})$, then divide.
\end{hint}

\begin{solution}
Two paths end in a positive test, so
\[
P(\text{pos}) = 0.0070 + 0.0990 = 0.1060,
\]
and
\[
P(H\mid\text{pos}) = \frac{P(H\cap\text{pos})}{P(\text{pos})}
= \frac{0.0070}{0.1060} \approx 0.0660,
\]
which is about $6.6\%$, or $7\%$ to the nearest percent.

This is the false-positive paradox again: heartworm is rare enough that the $9.9\%$
of healthy dogs who test positive vastly outnumber the $0.7\%$ of dogs that are both
sick and correctly flagged. A positive result on this test still means the dog
probably does not have heartworm --- which is why a positive screen is followed up
with a more specific confirmatory test.
\end{solution}
\end{prompt}
\end{problem}

\end{document}
